% !PRINTED 495
\weberpage{495}
%% [page furniture, not body text -- source-confirmed 2026-08-10]
%% \S\ 133. \hfill Die Resolventen 7ten und 11ten Grades.

%<<CONT
$A_\nu$ bleibt ungeändert durch die gleichzeitige Vertauschung:
\[
\left( f(\omega)^{12}, \frac{64}{f(\omega)^{12}} \right), \quad
\left( i\sqrt{11}, -i\sqrt{11} \right),
\]
und enthält bei geradem $\nu$ nur die geraden, bei ungeradem $\nu$
nur die ungeraden Potenzen von $f(\omega)^{12}$.

Auch hier treten nur Potenzen von $f(\omega)$ im Nenner auf, und die
Potenzen von $f(\omega)$ steigen bis zur selben Höhe wie die von
$f(\omega)^{-1}$.

Um über die Grade der Functionen $A_\nu$ ins Klare zu kommen,
betrachten wir die Anfänge der Entwickelungen nach steigenden
Potenzen von $q = e^{\pi i \omega}$.

Es beginnt die Entwickelung von $f(\omega)$ mit $q^{-\frac{1}{24}}$,
die von $v_\infty$ mit $q^{-\frac{n}{24}}$, von $v_c$ mit
$q^{-\frac{1}{24n}} e^{-\frac{2\pi i c}{48}\, \frac{1}{n}}$, worin
aber in der letzten Exponentialgrösse, wenn wir $c$ auf seinen
kleinsten Rest $(\mathrm{mod}\ n)$ reduciren, auch $\frac{1}{48}$
nach dem Modul $n$ zu nehmen, also für $n = 7, 11$:
\[
\frac{1}{48} \equiv -1 \ (\mathrm{mod}\ 7), \qquad
\frac{1}{48} \equiv 3 \ (\mathrm{mod}\ 11)
\]
zu setzen ist. Wenn wir also hiernach
\[
\varrho = e^{\frac{2\pi i}{7}} \quad \text{oder} \quad = e^{-\frac{6\pi i}{11}}
\]
setzen, so erhalten wir aus (14), (15) als Anfänge der Entwickelung
für $n=7$:
\[
i\sqrt{7}\, w_\nu = q^{-\frac{1}{7}} \varrho^{3\nu} (\varrho - \varrho^3)(\varrho^2 - \varrho^6)(\varrho^4 - \varrho^5) + \cdots,
\]
für $n=11$:
\[
i\sqrt{11}\, w_\nu = q^{-\frac{5}{22}} \varrho^{5\nu} (\varrho - \varrho^2)(\varrho^3 - \varrho^6)(\varrho^4 - \varrho^8)(\varrho^5 - \varrho^{10})(\varrho^9 - \varrho^7) + \cdots
\]
Es ist aber für $n = 7$:
\[
(\varrho - \varrho^3)(\varrho^2 - \varrho^6)(\varrho^4 - \varrho^5) = -\varrho - \varrho^2 - \varrho^4 + \varrho^3 + \varrho^5 + \varrho^6,
\]
und für $n = 11$:
\[
(\varrho - \varrho^2)(\varrho^3 - \varrho^6)(\varrho^4 - \varrho^8)(\varrho^5 - \varrho^{10})(\varrho^9 - \varrho^7) =
\]
\[
-\varrho - \varrho^3 - \varrho^4 - \varrho^5 - \varrho^9 +.\varrho^2 + \varrho^6 + \varrho^8 + \varrho^7 + \varrho^{10}. % [sic] stray period before $\varrho^2$ as printed
\]
